\begin{frame}{Top 10 de actividades económicas con mayores beneficios en impuesto de renta (2021)}
\begin{table}[h!]
    \centering
    \begin{tabular}{|c|l|p{3.5cm}|}
        \hline
        & \textbf{Actividad económica} & \textbf{Gasto tributario} (\% PIB) \\
        \hline
        1  & Transmisión de energía eléctrica & 0,13\% \\
        2  & Actividades de planes de seguridad social obligatoria & 0,11\% \\
        3  & Otras actividades del mercado de valores & 0,11\% \\
        4  & Banco central y Bancos comerciales & 0,09\% \\
        5  & Distribución de energía eléctrica & 0,07\% \\
        6  & Extracción de petróleo crudo & 0,06\% \\
        7  & Actividades de administración empresarial & 0,06\% \\
        8  & Seguros de vida & 0,05\% \\
        9  & Producción de malta y cervezas & 0,03\% \\
        10 & Otras actividades de servicio financiero & 0,03\% \\
        \hline
        \textbf{} & \textbf{Total} & \textbf{0,63\%} \\  
        \hline
    \end{tabular}
    \vspace{0.2cm}
    \par\scriptsize\textit{Fuente: Pardo (2024). El MFMP 2026 no actualiza este desglose sectorial; el último corte disponible en la fuente es 2021.}
\end{table}
\end{frame}
